\par
\textbf{Pour chaque affirmation, répondre par VRAI ou FAUX, en justifiant la réponse.}\par
Toute absence de justification ou justification incorrecte ne sera pas prise en compte dans la notation.

Une enquête est réalisée sur l'équipement informatique des 400 élèves d'un collège.
\begin{center}\renewcommand{\arraystretch}{1.35}
\begin{tabular}{|p{3.2cm}|c|c|c|}\hline
& Ont une adresse mail & N'ont pas d'adresse mail & Total\\\hline
Ont un équipement individuel & 50 & 210 & 260\\\hline
N'ont pas d'équipement individuel & 40 & 100 & 140\\\hline
Total & 90 & 310 & 400\\\hline
\end{tabular}
\end{center}

\begin{enumerate}
\item \underline{Affirmation}

50\% des élèves du collège possèdent une adresse mail et un équipement individuel.

\item \underline{Affirmation}

Au moins 50\% des élèves du collège ne possèdent pas d'adresse mail.

\item \underline{Affirmation}

On choisit un élève au hasard dans le collège. La probabilité qu'il ne possède ni adresse mail ni équipement individuel est égale à 25\%.

\item \underline{Affirmation}

Au moins $\dfrac{1}{5}$ des élèves ayant un équipement individuel possèdent également une adresse mail.
\end{enumerate}

\begin{color}{blue}Solution :\par

\begin{enumerate}
\item
50 élèves du collège possèdent une adresse mail et un équipement individuel sur 400 élèves au total. Cela correspond à une proportion de $\dfrac{50}{400} = \dfrac{12,5}{100} = 12,5~\% \neq 50~\% $.

\textbf{L'affirmation est fausse.}

\item 
310 élèves du collège ne possèdent pas d'adresse mail. $310>200$, la moitié de 400.

\textbf{L'affirmation est vraie.}
\item 

100 élèves du collège ne possèdent ni d'adresse mail ni d'équipement individuel sur 400 élèves au total. Cela correspond à une proportion de $\dfrac{100}{400} = \dfrac{25}{100} = 25~\% $. Si on choisit un élève au hasard parmi les élèves du collège, avec l'hypothèse d'équiprobabilité, la probabilité qu'il ne possède ni adresse mail ni équipement individuel est égale à 25\%.

\textbf{L'affirmation est vraie.}
\item 

260 élèves possèdent un équipement individuel. Parmi ceux-ci, 50 ont une adresse mail. La proportion d'élèves ayant une adresse mail parmi les élèves ayant un équipement individuel est $\dfrac{50}{260}=\dfrac{5}{26}$ et $\dfrac{5}{26}<\dfrac{5}{25}$. Comme $\dfrac{5}{25}=\dfrac{1}{5}$, $\dfrac{50}{260}<\dfrac{1}{5}$.

\textbf{L'affirmation est fausse.}

\end{enumerate}

\end{color}
